\documentclass[10pt,a4paper]{amsart}
\usepackage[margin=19mm]{geometry}
\usepackage{amsmath,amssymb,mathtools}
\usepackage[scr=rsfs,cal=euler]{mathalfa}
\usepackage{xfrac}
\usepackage[unicode,hidelinks,bookmarks=false]{hyperref}
\newcommand{\ie}{i.e.\ }
\newcommand{\cf}{cf.\ }
\newcommand{\resp}{respectively\ }
\newcommand{\Subsection}{\subsection}
\providecommand{\nobreakdash}{\nobreak\hbox{-}}
\setlength{\emergencystretch}{2em}
\multlinegap0pt
\usepackage{bm}
\usepackage{bbm}
\usepackage{dsfont}
\usepackage{xspace}
\usepackage{cancel}
\usepackage{mathtools}
\usepackage{amsthm}
\makeatletter
\@ifundefined{theorem}{\newtheorem{theorem}{Theorem}}{}
\@ifundefined{lemma}{\newtheorem{lemma}[theorem]{Lemma}}{}
\makeatother
\title[A Lower Bound on the Blob]{A Lower Bound on the Blob}
\author{James Schmidt}
\date{Source: arXiv:2609.13313v1 (CC BY 4.0)}
\begin{document}
\maketitle

\noindent\emph{Source and licence.} A Lower Bound on the Blob. Version arXiv:2609.13313v1 (CC BY 4.0), distributed under the Creative Commons Attribution 4.0 International licence, as recorded in the arXiv licence metadata for this exact version and shown on the abstract page https://arxiv.org/abs/2609.13313v1. Full rights notice: licenses/arxiv-2609.13313-CC-BY-4.0.txt.\par\smallskip

\noindent\textit{Editorial scope.} Counted unit: the Lemma whose source label is \texttt{lma:wall-reqs-rays} in the article (source0.tex lines 412--415, source label \texttt{lma:wall-reqs-rays}), delivered together with the article's own complete original proof of it (source0.tex lines 416--434, one \texttt{proof} environment of 19 lines). No mathematics is written, repaired or re-derived: statement and proof are the article's own text line for line. Required context retained uncounted: lma:front-erasure (lines 299--302); lma:creeping-wall (lines 356--359); lma:wall-surface-cont (lines 376--379); lma:loop-creeping (lines 387--390). Every definition, display and label of those lines is retained unchanged. Omitted from this package and not redistributed: 1--298, 303--355, 360--375, 380--386, 391--411, 435--748. Nothing inside a retained excerpt is omitted or altered: every retained body is delivered exactly as the article's lines are written, so no omission receipt of any kind accompanies this package. Citations: the retained excerpts carry no citation command at all, so no bibliography entry of the article is transcribed and no reference is redistributed. Reference adaptation: none was needed and none was made; every reference inside the delivered unit resolves inside it, so no unresolved reference or citation is produced. Printed slips kept literal: no printed word, symbol, hypothesis or number of the counted statement or of its proof was altered. Layout and class: the article's own class is \texttt{article} with packages including geometry, amsmath, amsthm, amssymb, amsfonts, graphicx, url, float; the wrapper is the lane's pinned \texttt{amsart} editorial wrapper at 10pt on A4, which restates the theorem-like environments the retained text needs and adds the article's own macro definitions that the retained text needs (none), with extra wrapper packages (bm, bbm, dsfont, xspace, cancel, mathtools). The delivered numbering is the unit's own. Assets: no retained span contains a figure environment, a graphics-inclusion command, a PDF-page inclusion, a table or a TikZ picture; no image file is copied into this package, the package declares no asset dependency and no image file is redistributed with it. Licence: version arXiv:2609.13313v1 of this article is distributed under the Creative Commons Attribution 4.0 International licence, as recorded in the arXiv licence metadata for this exact version and shown on the abstract page https://arxiv.org/abs/2609.13313v1. A full rights notice is delivered at licenses/arxiv-2609.13313-CC-BY-4.0.txt. Characters: every retained excerpt is pure ASCII, so the delivered engine, XeLaTeX with its default Latin Modern text font, reports no missing character.
\begin{lemma}[Front Erasure]
\label{lma:front-erasure}
Each front can only cease to exist when it runs into another front.
\end{lemma}

\begin{lemma}[Creeping Wall Requirement]
\label{lma:creeping-wall}
If a front creeps for \(x\) seconds, it touches a length of \(x\) surface. 
\end{lemma}

\begin{lemma}[Wall Surface Continuity]
\label{lma:wall-surface-cont}
If the vine touches a surface at some time \(t\), if the vine is allowed to creep forever with no more walls ever added, the vine will touch the opposing surface of the wall if and only if the wall whose surface the vine is touching is not part of a loop whose other face contains no point of the vine. 
\end{lemma}

\begin{lemma}[Loop Creeping]
\label{lma:loop-creeping}
If the vine touches a loop of wall, where the face on the other side of the loop there is no vine, the vine is either already successfully contained or will, if the vine is allowed to creep forever with no more walls ever added, creep around the loop and contain the entire loop strictly inside itself.
\end{lemma}

\begin{lemma}[Wall Requirements for Creeping on Rays]
\label{lma:wall-reqs-rays}
If, in the cardinal \(n\)-vine, a tendril moving along its ray touches a piece of wall and is halted at a time \(t\), and the tendril is un-halted at time \(t+x\) (and not before) to continue along its ray, then over the course of that time \(x\) there must have been at least two fronts creeping, each covering a distance of \(x\). 
\end{lemma}

\begin{proof}[Proof]
This is trivially true if \(x=0\). So assume \(x>0\). 

By Lemma~\ref{lma:creeping-wall}, we must only show that there are at least two fronts creeping between times \(t\) and \(t+x\). We may say that the wall along the ray which is reached by the vine at time \(t+x\) is, at time \(t\), untouched (if it is not yet built, then of course it must be untouched). Therefore, as the vine reaches points either by sending a tendril along that ray or by creeping, and the tendril along this ray is halted, the vine must reach this point by creeping. 

Therefore, at time \(t+x\), creeping must occur and thus there must be at least one front. 

By Lemma~\ref{lma:front-erasure}, fronts can only cease to exist when they hit another front, and as fronts are created in pairs when a tendril touches an untouched surface of wall or a new wall is drawn to touch a touched surface of wall, there must be an even number of fronts at any given time. 

Therefore, it suffices to show that, at no time between \(t\) and \(t+x\), were there zero fronts. 

There are zero fronts at a given time \(t\) if and only if there are no sections of surfaces touched by the vine adjacent to sections of surfaces untouched by the vine. In other words, every surface that is touched by the vine is only adjacent to surfaces that are also touched by the vine. As such, the whole connected curve of every surface which has any point of the vine must be entirely contained within the vine. 

By Lemma~\ref{lma:wall-surface-cont} and Lemma~\ref{lma:loop-creeping}, this would mean that either the vine has crept over to the other side of the wall that the tendril touched and was halted for, meaning the tendril is continuing along its ray, or the vine-touched side is contained within a loop. 

But then, either there is no vine outside of the loop, in which case it cannot possibly creep to the outside of the loop, or there is vine outside of the loop. If there is vine outside of the loop, then, as there is vine inside of the loop, there must be vine on both sides of the loop. As the vine is contiguous, there must be vine touching both the inside and the outside surfaces of the loop. Therefore, the entire loop must be touched by the vine inside and outside, meaning that the other side of the vine-touched side that halted the tendril is, in fact, vine-touched and thus the tendril is continuing along its ray already. 

Therefore, if there is no front active at any time \(t+z\) strictly between \(t\) and \(t+x\), then at that time the tendril will be un-halted and will be continuing along its ray at time \(t+z\). This is a contradiction. Therefore, at all times between \(t\) and \(t+x\), there are at least two fronts active. 
\end{proof}

\end{document}
